
retrouve la notion «~géométrique~» classique d'équivalence algébrique, envisagée dans 2.4.1. Aussi nous préférons dire qu'un élément de @@M0@@ est algébriquement équivalent à zéro si son image dans @@M1@@ l'est, où @@M2@@ est une clôture algébrique de @@M3@@, et de procéder de même pour la définition de l'équivalence algébrique dans les groupes @@M4@@, @@M5@@ et @@M6@@.

\textbf{7.12.} Ces définitions étant posées, je dis que \emph{l'homomorphisme de spécialisation} (7.3.1) \emph{transforme éléments algébriquement équivalents à zéro en éléments algébriquement équivalents à zéro}, et définit par suite par passage au quotient un homomorphisme
@@M7@@

Pour le voir, on se ramène immédiatement au cas où @@M8@@ est complet et @@M9@@ algébriquement clos (moyennant une compatibilité facile explicitée dans 7.15). Il suffit de voir que @@M10@@ transforme un élément de la forme (7.11.1) de @@M11@@ en un élément algébriquement équivalent à zéro dans @@M12@@. Ici @@M13@@, et nous écrirons en conséquence @@M14@@, @@M15@@ au lieu de @@M16@@, @@M17@@ comme plus haut. On sait que @@M18@@ est projective sur @@M19@@, donc (EGA IV 2.8.5) il existe un schéma projectif et plat @@M20@@ sur @@M21@@ dont la fibre générique soit @@M22@@. Utilisant la résolution des singularités des surfaces arithmétiques de Abhyankar [6], on voit qu'on peut même choisir pour @@M23@@ un schéma régulier. (N.B. La platitude est conservée par résolution, grâce à EGA IV 14.3.8). Comme @@M24@@ est propre sur @@M25@@, les points @@M26@@ et @@M27@@ de @@M28@@ rationnels sur @@M29@@ se prolongent en des sections @@M30@@, @@M31@@ de @@M32@@ sur @@M33@@, et comme @@M34@@ est régulier, il est bien connu que @@M35@@ est lisse sur @@M36@@ en les points de @@M37@@, @@M38@@ (EGA IV, 17.12.1 c) @@M39@@ a) et 19.1.1). Par
